\documentclass[10pt,a4paper]{amsart}
\usepackage{amsmath,amssymb,amsthm,graphicx}
\usepackage{verbatim}
\newenvironment{reference}{\comment}{\endcomment}
%\newenvironment{reference}{}{}
\newenvironment{slogan}{\comment}{\endcomment}
\newenvironment{history}{\comment}{\endcomment}

% For commutative diagrams we use Xy-pic
\usepackage[all]{xy}

% We use 2cell for 2-commutative diagrams.
\xyoption{2cell}
% Theorem environments.
%
\theoremstyle{plain}
\newtheorem{theorem}[subsection]{Theorem}
\newtheorem{proposition}[subsection]{Proposition}
\newtheorem{lemma}[subsection]{Lemma}

\theoremstyle{definition}
\newtheorem{definition}[subsection]{Definition}
\newtheorem{example}[subsection]{Example}
\newtheorem{exercise}[subsection]{Exercise}
\newtheorem{situation}[subsection]{Situation}

\theoremstyle{remark}
\newtheorem{remark}[subsection]{Remark}
\newtheorem{remarks}[subsection]{Remarks}

\numberwithin{equation}{subsection}

% Macros
%
\def\lim{\mathop{\mathrm{lim}}\nolimits}
\def\colim{\mathop{\mathrm{colim}}\nolimits}
\def\Spec{\mathop{\mathrm{Spec}}}
\def\Hom{\mathop{\mathrm{Hom}}\nolimits}
\def\Ext{\mathop{\mathrm{Ext}}\nolimits}
\def\SheafHom{\mathop{\mathcal{H}\!\mathit{om}}\nolimits}
\def\SheafExt{\mathop{\mathcal{E}\!\mathit{xt}}\nolimits}
\def\Sch{\mathit{Sch}}
\def\Mor{\mathop{\mathrm{Mor}}\nolimits}
\def\Ob{\mathop{\mathrm{Ob}}\nolimits}
\def\Sh{\mathop{\mathit{Sh}}\nolimits}
\def\NL{\mathop{N\!L}\nolimits}
\def\CH{\mathop{\mathrm{CH}}\nolimits}
\def\proetale{{pro\text{-}\acute{e}tale}}
\def\etale{{\acute{e}tale}}
\def\QCoh{\mathit{QCoh}}
\def\Ker{\mathop{\mathrm{Ker}}}
\def\Im{\mathop{\mathrm{Im}}}
\def\Coker{\mathop{\mathrm{Coker}}}
\def\Coim{\mathop{\mathrm{Coim}}}

% Boxtimes
%
\DeclareMathSymbol{\boxtimes}{\mathbin}{AMSa}{"02}

%
% Macros for moduli stacks/spaces
%
\def\QCohstack{\mathcal{QC}\!\mathit{oh}}
\def\Cohstack{\mathcal{C}\!\mathit{oh}}
\def\Spacesstack{\mathcal{S}\!\mathit{paces}}
\def\Quotfunctor{\mathrm{Quot}}
\def\Hilbfunctor{\mathrm{Hilb}}
\def\Curvesstack{\mathcal{C}\!\mathit{urves}}
\def\Polarizedstack{\mathcal{P}\!\mathit{olarized}}
\def\Complexesstack{\mathcal{C}\!\mathit{omplexes}}
% \Pic is the operator that assigns to X its picard group, usage \Pic(X)
% \Picardstack_{X/B} denotes the Picard stack of X over B
% \Picardfunctor_{X/B} denotes the Picard functor of X over B
\def\Pic{\mathop{\mathrm{Pic}}\nolimits}
\def\Picardstack{\mathcal{P}\!\mathit{ic}}
\def\Picardfunctor{\mathrm{Pic}}
\def\Deformationcategory{\mathcal{D}\!\mathit{ef}}
\usepackage{cleveref}
\crefname{originallemma}{original lemma}{original lemmas}
\Crefname{originallemma}{Original lemma}{Original lemmas}
\makeatletter
\newlabel{lemma-adjoint}{{0778}{source}{}{sourcepointer-1}{}}
\newlabel{lemma-adjoint@cref}{{[originallemma][0][]0778}{[1][0][]source}{}{}{}}
\makeatother
\crefname{originalreference}{original reference}{original references}
\Crefname{originalreference}{Original reference}{Original references}
\makeatletter
\newlabel{categories-lemma-colimits-coproducts-coequalizers}{{002P}{source}{}{sourcepointer-2}{}}
\newlabel{categories-lemma-colimits-coproducts-coequalizers@cref}{{[originalreference][0][]002P}{[1][0][]source}{}{}{}}
\makeatother
\title{Kernel of the Q functor is the parasitic subcategory; Q is exact and bicontinuous}
\author{The Stacks Project Authors}
\date{stacks-project:stacks-cohomology.tex:local-snapshot-20261012; GNU FDL 1.2 or later}
\begin{document}
\maketitle
\noindent\textit{Editorial note.} All delivered mathematical text is the author's own, copied verbatim from the pinned Stacks Project chapter file stacks-cohomology.tex; no statement or proof step was rewritten, reordered or removed.\par\medskip
\noindent\textit{Editorial note.} The counted span cites the lemmas recorded in assembly.external\_references with canonical tags as navigation pointers.\par\medskip
\noindent\textit{Allowed standard prerequisites (editorial).}\begin{itemize}
\item Stack presentation by a smooth groupoid. Algebraic Stacks, Lemma algebraic-lemma-stack-presentation. Provides X = [U/R] with U a scheme.
\item Colimits from coproducts and coequalizers. Categories, Lemma categories-lemma-colimits-coproducts-coequalizers. Reduces colimit commutation to direct sums for an exact functor.
\item Parasitic modules form a Serre subcategory. Cohomology of Algebraic Stacks, Lemma stacks-cohomology-lemma-parasitic (2)(b). Named inside the proof for the snake-lemma step.
\end{itemize}\medskip
\ifdefined\hypertarget\hypertarget{sourcepointer-1}{}\fi
\noindent\textit{Source reference.} Original lemma 0778: sources/stacks-chow/tags.txt, source lines 9193--9193. Cross-chapter standard prerequisite cited by the author inside the counted span; it is used as a pointer and its own proof is not part of this item.\par\medskip
\ifdefined\hypertarget\hypertarget{sourcepointer-2}{}\fi
\noindent\textit{Source reference.} Original reference 002P: sources/stacks-chow/tags.txt, source lines 106--106. Cross-chapter standard prerequisite cited by the author inside the counted span; it is used as a pointer and its own proof is not part of this item.\par\medskip
\begin{lemma}
\label{lemma-adjoint-kernel-parasitic}
Let $\mathcal{X}$ be an algebraic stack.
Let $Q : \textit{LQCoh}^{fbc}(\mathcal{O}_\mathcal{X}) \to
\QCoh(\mathcal{O}_\mathcal{X})$
be the functor constructed in Lemma \ref{lemma-adjoint}.
\begin{enumerate}
\item The kernel of $Q$ is exactly the collection of parasitic objects
of $\textit{LQCoh}^{fbc}(\mathcal{O}_\mathcal{X})$.
\item For any object $\mathcal{F}$
of $\textit{LQCoh}^{fbc}(\mathcal{O}_\mathcal{X})$ both the kernel and the
cokernel of the
adjunction map $Q(\mathcal{F}) \to \mathcal{F}$ are parasitic.
\item The functor $Q$ is exact and commutes with all limits and colimits.
\end{enumerate}
\end{lemma}

\begin{proof}
Write $\mathcal{X} = [U/R]$ as in the proof of Lemma \ref{lemma-adjoint}.
Let $\mathcal{F}$ be an object of
$\textit{LQCoh}^{fbc}(\mathcal{O}_\mathcal{X})$.
It is clear from the proof of Lemma \ref{lemma-adjoint}
that $\mathcal{F}$ is in the kernel of $Q$ if and only if
$\mathcal{F}|_{U_\etale} = 0$.
In particular, if $\mathcal{F}$ is parasitic then $\mathcal{F}$ is in
the kernel. Next, let $x : V \to \mathcal{X}$ be a flat morphism, where
$V$ is a scheme. Set $W = V \times_\mathcal{X} U$ and consider the diagram
$$
\xymatrix{
W \ar[d]_p \ar[r]_q & V \ar[d] \\
U \ar[r] & \mathcal{X}
}
$$
Note that the projection $p : W \to U$ is flat and the projection
$q : W \to V$ is smooth and surjective. This implies that $q_{small}^*$
is a faithful functor on quasi-coherent modules. By assumption $\mathcal{F}$
has the flat base change property so that we obtain
$p_{small}^*\mathcal{F}|_{U_\etale} \cong
q_{small}^*\mathcal{F}|_{V_\etale}$. Thus if $\mathcal{F}$
is in the kernel of $Q$, then $\mathcal{F}|_{V_\etale} = 0$
which completes the proof of (1).

\medskip\noindent
Part (2) follows from the discussion above and the fact
that the map $Q(\mathcal{F}) \to \mathcal{F}$ becomes an isomorphism after
restricting to $U_\etale$.

\medskip\noindent
To see part (3) note that $Q$ is left exact as a right adjoint.
Let $0 \to \mathcal{F} \to \mathcal{G} \to \mathcal{H} \to 0$
be a short exact sequence in $\textit{LQCoh}^{fbc}(\mathcal{O}_\mathcal{X})$.
Consider the following commutative diagram
$$
\xymatrix{
0 \ar[r] &
Q(\mathcal{F}) \ar[r] \ar[d]_a &
Q(\mathcal{G}) \ar[r] \ar[d]_b &
Q(\mathcal{H}) \ar[r] \ar[d]_c & 0 \\
0 \ar[r] &
\mathcal{F} \ar[r] &
\mathcal{G} \ar[r] &
\mathcal{H} \ar[r] & 0
}
$$
Since the kernels and cokernels of $a$, $b$, and $c$ are parasitic by
part (2) and since the bottom row is a short exact sequence, we see that
the top row as a complex of $\mathcal{O}_\mathcal{X}$-modules has parasitic
cohomology sheaves (details omitted; this uses that the category of
parasitic modules is a Serre subcategory of the category of all modules).
By left exactness of $Q$ we see that only exactness at $Q(\mathcal{H})$
is at issue. However, the cokernel $\mathcal{Q}$ of
$Q(\mathcal{G}) \to Q(\mathcal{H}))$
may be computed either in $\textit{Mod}(\mathcal{O}_\mathcal{X})$
or in $\QCoh(\mathcal{O}_\mathcal{X})$ with the same result because
the inclusion functor $\QCoh(\mathcal{O}_\mathcal{X}) \to
\textit{LQCoh}^{fbc}(\mathcal{O}_\mathcal{X})$
is a left adjoint and hence right exact. Hence $\mathcal{Q} = Q(\mathcal{Q})$
is both quasi-coherent and parasitic, whence $0$ by part (1) as desired.

\medskip\noindent
As a right adjoint $Q$ commutes with all limits. Since $Q$ is exact, to
show that $Q$ commutes with all colimits it suffices to show that $Q$
commutes with direct sums, see
Categories, Lemma \ref{categories-lemma-colimits-coproducts-coequalizers}.
Let $\mathcal{F}_i$, $i \in I$ be a family of objects of
$\textit{LQCoh}^{fbc}(\mathcal{O}_\mathcal{X})$. To see that
$Q(\bigoplus \mathcal{F}_i)$ is equal to $\bigoplus Q(\mathcal{F}_i)$
we look at the construction of $Q$ in the proof of
Lemma \ref{lemma-adjoint}. This uses a presentation $\mathcal{X} = [U/R]$
where $U$ is a scheme. Then $Q(\mathcal{F})$
is computed by first taking the pair
$(\mathcal{F}|_{U_\etale}, \alpha)$ in $\QCoh(U, R, s, t, c)$
and then using the equivalence
$\QCoh(U, R, s, t, c) \cong \QCoh(\mathcal{O}_\mathcal{X})$.
Since the restriction functor $\textit{Mod}(\mathcal{O}_\mathcal{X}) \to
\textit{Mod}(\mathcal{O}_{U_\etale})$,
$\mathcal{F} \mapsto \mathcal{F}|_{U_\etale}$ commutes with
direct sums, the desired equality is clear.
\end{proof}

\end{document}
